\section{Sujet n\textsuperscript{o}\,5}
\noindent\begin{minipage}{0.5\linewidth}\footnotesize Office du Baccalauréat du Cameroun\end{minipage}%
\hfill\begin{minipage}{0.45\linewidth}\footnotesize\raggedleft Examen : \textbf{Baccalauréat} · Série : \textbf{C/E}\\
Épreuve : \textbf{Mathématiques} · Durée : \textbf{4\,h} · Coef. : \textbf{7}\end{minipage}
\par\smallskip{\color{onbo}\rule{\linewidth}{1pt}}

\subsection*{Partie A — Évaluation des ressources \hfill (15 points)}

\noindent\textbf{Exercice 1.} \hfill\textit{(3 points)}\\
Une coopérative de transport relève, sur $7$ de ses lignes, la distance $x$ (en \si{\kilo\metre})
et la recette journalière $y$ (en milliers de FCFA) :
\begin{center}\small
\begin{tabular}{c|ccccccc}
\toprule
$x$ & 4 & 6 & 8 & 10 & 12 & 14 & 16\\
$y$ & 9 & 12 & 14 & 19 & 21 & 26 & 27\\
\bottomrule
\end{tabular}
\end{center}
\begin{enumerate}[label=\arabic*)]
\item Calculer les moyennes $\bar x$ et $\bar y$. \hfill\textit{(0,75 pt)}
\item Calculer le coefficient de corrélation linéaire $r$ et interpréter. On donne
$\sigma_x^2=16$, $\sigma_y^2\approx41{,}43$ et $\mathrm{cov}(x,y)=25$. \hfill\textit{(1,25 pt)}
\item Donner une équation de la droite de régression de $y$ en $x$ par la méthode des moindres carrés. \hfill\textit{(1 pt)}
\end{enumerate}

\smallskip
\noindent\textbf{Exercice 2.} \hfill\textit{(3 points)}\\
On modélise le nombre de véhicules de la coopérative par la suite $(u_n)$ définie par
$u_0=20$ et $u_{n+1}=0{,}8\,u_n+10$ pour tout $n\in\N$.
\begin{enumerate}[label=\arabic*)]
\item On pose $v_n=u_n-50$. Montrer que $(v_n)$ est géométrique ; préciser sa raison et $v_0$. \hfill\textit{(1,5 pt)}
\item Exprimer $u_n$ en fonction de $n$, puis déterminer la limite de $(u_n)$ et l'interpréter. \hfill\textit{(1,5 pt)}
\end{enumerate}

\smallskip
\noindent\textbf{Exercice 3.} \hfill\textit{(4 points)}\\
Soit $f$ la fonction définie sur $\R$ par $f(x)=(2-x)\mathrm{e}^{x}$, de courbe $(\mathcal{C})$.
\begin{enumerate}[label=\arabic*)]
\item Déterminer les limites de $f$ en $-\infty$ et $+\infty$. \hfill\textit{(1 pt)}
\item Montrer que $f'(x)=(1-x)\mathrm{e}^{x}$ et dresser le tableau de variations de $f$. \hfill\textit{(1,5 pt)}
\item À l'aide d'une intégration par parties, calculer $\displaystyle J=\int_0^{1}(2-x)\mathrm{e}^{x}\,\dd x$. \hfill\textit{(1,5 pt)}
\end{enumerate}

\smallskip
\noindent\textbf{Exercice 4.} \hfill\textit{(5 points)}\\
Le plan complexe est muni d'un repère orthonormé direct $(O\,;\,\vec u,\vec v)$. On considère
les points $A$, $B$, $C$ d'affixes respectives $z_A=1+i$, $z_B=4+i$ et $z_C=1+5i$.
\begin{enumerate}[label=\arabic*)]
\item Calculer $\dfrac{z_C-z_A}{z_B-z_A}$ et en déduire la nature du triangle $ABC$. \hfill\textit{(1,5 pt)}
\item Soit $S$ la similitude directe qui transforme $A$ en $B$ et $B$ en $C$.
Déterminer l'écriture complexe $z'=az+b$ de $S$. \hfill\textit{(1,5 pt)}
\item Déterminer le rapport et l'angle de $S$. \hfill\textit{(1 pt)}
\item Déterminer l'affixe $\omega$ du centre $\Omega$ de $S$. \hfill\textit{(1 pt)}
\end{enumerate}

\subsection*{Partie B — Évaluation des compétences \hfill (5 points)}

\noindent\textit{La coopérative «~La Confiance~» gère des taxis et des bus à Douala. Elle te
sollicite, en tant qu'élève de Terminale~C, pour trois décisions. Elle veut d'abord aménager
un parking rectangulaire \emph{le long d'un mur} (le côté longeant le mur ne se barrière pas)
avec \SI{160}{\metre} de barrière. Son chiffre d'affaires, de $30$ millions FCFA cette première
année, augmente de \SI{12}{\percent} par an. Enfin, sur sa flotte, \SI{4}{\percent} des véhicules
tombent en panne un jour donné ; un inspecteur en contrôle $6$ au hasard.}

\begin{enumerate}[label=\textbf{Tâche \arabic*.},leftmargin=2.4em]
\item Déterminer les dimensions du parking rectangulaire d'\emph{aire maximale} qu'il peut
clôturer, ainsi que cette aire. \hfill\textit{(1,5 pt)}
\item Déterminer l'année à partir de laquelle son chiffre d'affaires dépassera $60$ millions FCFA. \hfill\textit{(1,5 pt)}
\item Déterminer la probabilité qu'\emph{au moins un} des $6$ véhicules contrôlés soit en panne. \hfill\textit{(1,5 pt)}
\end{enumerate}
\par\smallskip{\footnotesize\itshape Présentation générale : 0,5 point.}

% ════════════════════════════════════════════════════
\corrige{du sujet 5}

\noindent\textbf{Partie A — Exercice 1.}
1) $\bar x=\frac{4+6+8+10+12+14+16}{7}=\frac{70}{7}=10$ ;
$\bar y=\frac{9+12+14+19+21+26+27}{7}=\frac{128}{7}\approx18{,}29$.
2) $r=\dfrac{\mathrm{cov}(x,y)}{\sigma_x\,\sigma_y}=\dfrac{25}{\sqrt{16}\times\sqrt{41{,}43}}
=\dfrac{25}{4\times6{,}437}\approx0{,}971$. Très proche de $1$ : forte corrélation linéaire positive.
3) Pente $a=\dfrac{\mathrm{cov}(x,y)}{\sigma_x^2}=\dfrac{25}{16}=1{,}5625$ ;
ordonnée $b=\bar y-a\bar x=18{,}29-15{,}625\approx2{,}66$.
Droite : $y=1{,}56\,x+2{,}66$.

\noindent\textbf{Exercice 2.}
1) $v_{n+1}=u_{n+1}-50=0{,}8u_n+10-50=0{,}8u_n-40=0{,}8(u_n-50)=0{,}8\,v_n$.
Donc $(v_n)$ est géométrique de raison $q=0{,}8$ et de premier terme $v_0=u_0-50=-30$.
2) $v_n=-30\times0{,}8^{\,n}$, d'où $u_n=50-30\times0{,}8^{\,n}$.
Comme $0<0{,}8<1$, $0{,}8^{\,n}\to0$, donc $u_n\to50$ : la flotte se stabilise à $50$ véhicules.

\noindent\textbf{Exercice 3.}
1) En $-\infty$ : $\mathrm{e}^x\to0$ et par croissances comparées $(2-x)\mathrm{e}^x\to0^+$ (donc $f\to0$).
En $+\infty$ : $2-x\to-\infty$ et $\mathrm{e}^x\to+\infty$, donc $f(x)\to-\infty$.
2) $f'(x)=(-1)\mathrm{e}^x+(2-x)\mathrm{e}^x=(1-x)\mathrm{e}^x$. $\mathrm{e}^x>0$, donc $f'$ a le signe de $1-x$ :
$f$ croît sur $]-\infty,1]$, décroît sur $[1,+\infty[$, maximum $f(1)=(2-1)\mathrm e=\mathrm e$.
3) IPP ($u=2-x$, $v'=\mathrm{e}^x$, $u'=-1$, $v=\mathrm e^x$) :
$J=\big[(2-x)\mathrm{e}^x\big]_0^1+\int_0^1\mathrm{e}^x\,\dd x=(\mathrm e-2)+\big[\mathrm e^x\big]_0^1
=(\mathrm e-2)+(\mathrm e-1)=2\mathrm e-3\approx2{,}44$.

\noindent\textbf{Exercice 4.}
1) $z_B-z_A=3$, $z_C-z_A=4i$, donc $\dfrac{z_C-z_A}{z_B-z_A}=\dfrac{4i}{3}=\dfrac43 i$.
Le quotient est imaginaire pur : $(AC)\perp(AB)$, donc le triangle $ABC$ est rectangle en $A$.
2) $S(A)=B$ et $S(B)=C$ : $a z_A+b=z_B$ et $a z_B+b=z_C$. Par soustraction
$a(z_B-z_A)=z_C-z_B$, soit $a=\dfrac{z_C-z_B}{z_B-z_A}=\dfrac{(1+5i)-(4+i)}{3}=\dfrac{-3+4i}{3}=-1+\tfrac43 i$.
Puis $b=z_B-a z_A=(4+i)-(-1+\tfrac43 i)(1+i)$. Or $(-1+\tfrac43 i)(1+i)=-1-i+\tfrac43 i+\tfrac43 i^2
=-1-i+\tfrac43 i-\tfrac43=-\tfrac73+\tfrac13 i$, d'où $b=(4+i)-(-\tfrac73+\tfrac13 i)=\tfrac{19}{3}+\tfrac23 i$.
Écriture : $z'=\big(-1+\tfrac43 i\big)z+\tfrac{19}{3}+\tfrac23 i$.
3) $|a|=\sqrt{1+\tfrac{16}{9}}=\sqrt{\tfrac{25}{9}}=\tfrac53$ : rapport $\tfrac53$.
$\arg a=\arg(-1+\tfrac43 i)=\pi-\arctan\tfrac{4}{3}\approx\pi-0{,}927\approx2{,}214$ rad
(soit environ $126{,}87^\circ$) : c'est l'angle de $S$.
4) Le centre vérifie $\omega=a\omega+b$, soit $\omega(1-a)=b$ avec $1-a=2-\tfrac43 i$.
$\omega=\dfrac{b}{1-a}=\dfrac{\tfrac{19}{3}+\tfrac23 i}{2-\tfrac43 i}=\dfrac{19+2i}{6-4i}
=\dfrac{(19+2i)(6+4i)}{36+16}=\dfrac{114+76i+12i+8i^2}{52}=\dfrac{106+88i}{52}=\dfrac{53}{26}+\dfrac{22}{13}i$.

\smallskip
\noindent\textbf{Partie B — Situation.}
\emph{Tâche 1.} On note $x$ la dimension des deux côtés perpendiculaires au mur et $y=160-2x$
celle parallèle au mur. Aire $A(x)=x(160-2x)$ ; $A'(x)=160-4x=0$ en $x=40$.
Dimensions : $x=\SI{40}{\metre}$, $y=\SI{80}{\metre}$ ; aire maximale $=40\times80=\SI{3200}{\metre\squared}$.

\emph{Tâche 2.} Chiffre d'affaires de l'année $n$ : $u_n=30\times1{,}12^{\,n-1}$. On résout
$u_n>60$, soit $1{,}12^{\,n-1}>2$, d'où $n-1>\dfrac{\ln 2}{\ln 1{,}12}\approx6{,}12$ : à partir de la
\textbf{8\textsuperscript{e} année} ($u_8=30\times1{,}12^{7}\approx66{,}3$ millions FCFA).

\emph{Tâche 3.} $P(\text{au moins un en panne})=1-(0{,}96)^6\approx1-0{,}783=\mathbf{0{,}217}$,
soit environ \SI{21,7}{\percent}.
